% STATUS: written 2026-11-05
% Sources read in full: Mikhalev-Pankrat'ev 1989 (J. Soviet Math. 45,
% 912--955, p. 912); Hardouin-Singer 2008 (arxiv version, Definitions
% 2.1, 2.5, Theorem 2.6); Pillay-Ziegler 2003 (Selecta Math. 9,
% 579--599, p. 579).
\section{Difference and differential--difference algebra}\label{sec:differential-algebra:difference}

So far this chapter has used one operator, the
derivation~\cite[p.~912]{MikhalevPankratev1989}. A small change replaces the
whole theory: use an automorphism $\sigma$ of the ring in place of the
derivation~\cite[p.~912]{MikhalevPankratev1989}. Ritt named the resulting
theory \emph{difference algebra}. The monograph of Cohn is its standard
reference~\cite[p.~912]{MikhalevPankratev1989}. Pillay and Ziegler describe
the change in geometric language, following
Section~\ref{sec:differential-algebra:model-theory}. Difference algebraic
geometry is the analogue of differential algebraic geometry ``when the
derivation is replaced by an
automorphism''~\cite[p.~579]{PillayZiegler2003}.

The two theories run in parallel~\cite[p.~912]{MikhalevPankratev1989}. Each
main result of differential algebra has a difference analogue, which appears a
little later~\cite[p.~912]{MikhalevPankratev1989}. The table gives the
dictionary of the two sides, as Mikhalev and Pankrat'ev set it
out~\cite[p.~912]{MikhalevPankratev1989}.

\begin{center}
\renewcommand{\arraystretch}{1.2}
\setlength{\tabcolsep}{4pt}
\begin{tabular}{ll}
\textbf{differential side} & \textbf{difference side} \\[2pt]
derivation $\partial$ & automorphism $\sigma$. \\
differential field & difference field. \\
differential ideal & difference ideal. \\
Ritt--Raudenbush basis & difference basis theorem. \\
Picard--Vessiot extension & Hardouin--Singer ring. \\
integration in finite terms & summation in finite terms.
\end{tabular}
\end{center}

Mikhalev and Pankrat'ev warn that the dictionary is not a
translation~\cite[p.~912]{MikhalevPankratev1989}. Their words are that ``all
these results in difference algebra are essentially different from the
corresponding results in differential
algebra''~\cite[p.~912]{MikhalevPankratev1989}. A $\sigma$-ideal need not be
radical when its differential analogue is, and $\sigma^{-1}$ plays a role the
derivation has no analogue of~\cite[p.~912]{MikhalevPankratev1989}. This
section reports the definitions and the Galois theorem on the difference side,
as Hardouin and Singer state them~\cite[pp.~5--6]{HardouinSinger2008}.

The clearest formulation uses both operators and lets them
commute~\cite[Definition~2.1, p.~5]{HardouinSinger2008}.

\begin{definition}[{\cite[Definition~2.1, p.~5]{HardouinSinger2008}}]\label{def:diffalg:sigma-delta-ring}
A \emph{$\sigma\partial$-ring} is a commutative ring $R$ with unit
together with an automorphism $\sigma$ and a derivation $\partial$
satisfying $\sigma(\partial(r)) = \partial(\sigma(r))$ for all
$r \in R$. A \emph{$\sigma\partial$-field} is a $\sigma\partial$-ring
which is a field.
\end{definition}

The three basic examples already appear in the source
\cite[Example~2.2, p.~5]{HardouinSinger2008}.

\begin{example}[{\cite[Example~2.2, p.~5]{HardouinSinger2008}}]\label{ex:diffalg:sigma-delta-examples}
The following are $\sigma\partial$-fields.
\begin{enumerate}
\item $\mathbb{C}(x)$ with $\sigma(x) = x+1$ and $\partial = d/dx$;
\item $\mathbb{C}(x)$ with $\sigma(x) = qx$ ($q\in\mathbb{C}\setminus\{0\}$)
  and $\partial = x\,d/dx$;
\item $\mathbb{C}(x,t)$ with $\sigma(x) = x+1$, $\sigma(t)=t$ and
  $\partial = \partial/\partial t$.
\end{enumerate}
\end{example}

The equation of interest is the difference equation $\sigma(Y) = AY$, taken
over a $\sigma\partial$-field $k$ with $A$ an invertible matrix over
$k$~\cite[pp.~5--6]{HardouinSinger2008}. To it one attaches a
$\sigma\partial$-\emph{Picard--Vessiot ring}. This is a ring containing $k$
with no proper $\sigma\partial$-ideal~\cite[Definition~2.3, p.~5]{HardouinSinger2008}.
The ring holds a fundamental solution
matrix~\cite[Definition~2.3, p.~5]{HardouinSinger2008}. It is generated over
$k$, as a
$\partial$-ring, by the entries of that matrix and by the inverse of its
determinant~\cite[Definition~2.3, p.~5]{HardouinSinger2008}. The
$\sigma\partial$-Galois group is the group of its
automorphisms~\cite[Definition~2.5, p.~6]{HardouinSinger2008}.

\begin{definition}[{\cite[Definition~2.5, p.~6]{HardouinSinger2008}}]\label{def:diffalg:sigma-delta-galois-group}
The \emph{$\sigma\partial$-Galois group} $\mathrm{Aut}_{\sigma\partial}(R/k)$
of a $\sigma\partial$-Picard--Vessiot ring $R$ over a
$\sigma\partial$-field $k$ is the group of all
$\sigma\partial$-$k$-automorphisms of $R$.
\end{definition}

When $\partial$ vanishes identically, this is the Galois group of the
difference equation alone, which is a linear algebraic
group~\cite[pp.~5--6]{HardouinSinger2008}. When $\partial$ is present, the
group carries an additional structure and lives among the linear differential
algebraic groups~\cite[Theorem~2.6, p.~6]{HardouinSinger2008}. That is the
content of the main theorem of Hardouin and
Singer~\cite[Theorem~2.6, p.~6]{HardouinSinger2008}.

\begin{theorem}[{\cite[Theorem~2.6, p.~6]{HardouinSinger2008}}]\label{thm:diffalg:sigma-delta-group}
Let $k$ be a $\sigma\partial$-field whose fixed field $k^\sigma$ is a
differentially closed field, and let $R = k\{Z, \det Z^{-1}\}_\partial$,
$\sigma(Z) = AZ$, be a $\sigma\partial$-Picard--Vessiot extension of $k$.
\begin{enumerate}
\item The Galois group $\mathrm{Aut}_{\sigma\partial}(R/k)$ identifies
  with the set of $k^\sigma$-points of a linear
  $\partial$-differential algebraic group $G$ defined over $k^\sigma$.
\item The ring $R$ is reduced and is the coordinate ring of a $G$-torsor
  $V$ defined over $k$; the action of $G$ on $V$ induces the action of
  $\mathrm{Aut}_{\sigma\partial}(R/k)$ on $R$.
\end{enumerate}
\end{theorem}

So $\sigma$ and $\partial$ together give back exactly the two theories that
they extend~\cite[pp.~5--6]{HardouinSinger2008}. Take $\sigma$ trivial, and
Theorem~\ref{thm:diffalg:sigma-delta-group} reduces to the Picard--Vessiot
Theorem~\ref{thm:diffalg:pv-fundamental}. Take $\partial$ trivial, and it
returns the algebraic-group Galois theory of a pure difference
equation~\cite[pp.~5--6]{HardouinSinger2008}.

The pure difference side has a recent theorem of its own
\cite[Theorem~1.2, p.~2]{WibmerBachmayr2025}. Wibmer and Bachmayr ask when
every torsor for a linear difference algebraic group is trivial after a finite
extension of the difference field~\cite[Theorem~1.2, p.~2]{WibmerBachmayr2025}.
They answer with three equivalent conditions, and the third names the
obstruction~\cite[Theorem~1.2, p.~2]{WibmerBachmayr2025}.

\begin{theorem}[{\cite[Theorem~1.2, p.~2]{WibmerBachmayr2025}}]\label{thm:diffalg:trivializing-torsors}
For an ordinary difference field $k$ of characteristic zero the following
statements are equivalent.
\begin{enumerate}
\item For every linear difference algebraic group $G$ over $k$ and every
$G$-torsor $X$, there is an integer $n \ge 1$ with $X(k^{[n]}) \neq \emptyset$.
\item The direct limit of the first difference cohomology sets
$H^1_{\sigma^n}(k, G_n)$ vanishes for every linear difference algebraic group
$G$ over $k$.
\item The field $k$ is algebraically closed, and the matrix equation
$\sigma(Y) = AY$ and the one-relator equation
$y^{\alpha_0}\sigma(y)^{\alpha_1}\cdots\sigma^m(y)^{\alpha_m} = a$ are each
solvable in some $k^{[n]}$.
\end{enumerate}
\end{theorem}

Wibmer and Bachmayr note what the theorem does not
assume~\cite[p.~3]{WibmerBachmayr2025}. No condition is placed on the constants
$k^\sigma$ of $k$~\cite[p.~3]{WibmerBachmayr2025}. The one-relator subgroups
of the multiplicative group are, in their words, the only obstruction to the
immediate difference analogue of the differential
statement~\cite[p.~3]{WibmerBachmayr2025}.

This is how the chapter ends: differential, model-theoretic, and difference
algebra are three expressions of one algebraic method, applied to three kinds
of operator~\cite[p.~912]{MikhalevPankratev1989}. The difference side also
carries its own recent results, of which
Theorem~\ref{thm:diffalg:trivializing-torsors} is one.

A Hilbert-type bivariate quasi-polynomial is attached to a difference-differential field extension with weighted basic derivations, and its existence and invariants are determined~\cite{WlBivariateDimensionQuasiPolynom2020}.
